% =====================================================================
\chapter{Proyecto de fase}
\label{cap:proyecto}
% =====================================================================

\section*{Informe técnico del modelo «Alcotán»}
\phantomsection\addcontentsline{toc}{section}{Informe técnico del modelo «Alcotán»}

\paragraph{Objetivo.} El estudiante actúa como equipo de sistemas de un laboratorio ficticio y debe proponer y justificar el modelo Alcotán: cuánto cómputo hay disponible, qué tamaño y cuántos datos usar, cómo repartir el entrenamiento en el clúster y cómo servir el modelo. Alcotán es un modelo inventado; ninguna de sus cifras debe atribuirse a un modelo real. El proyecto integra los cuatro módulos de la fase: leyes de escalado (M16), sistemas (M17), arquitectura (M18) e inferencia (M19).

\paragraph{Restricciones.}
\begin{itemize}[label=\textendash]
  \item Clúster: 2048 H100 (256 nodos de 8 GPU, NVLink dentro del nodo e InfiniBand entre nodos) durante 60 días, con una MFU objetivo del 40\,\% (cifras aproximadas).
  \item Electricidad: 700~W por GPU y $\PUE=1{,}2$. Precio de alquiler: 2{,}5~\$ por GPU-hora.
  \item Servicio: al menos 30 tokens/s por usuario con 64 usuarios simultáneos y 4096 tokens de contexto, en un nodo de 8 H100.
  \item Datos: $1{,}2\cdot 10^{13}$ tokens únicos tras filtrar y deduplicar.
\end{itemize}

\paragraph{Entregables.}
\begin{enumerate}[label=\arabic*.]
  \item Informe de 6 a 10 páginas (PDF o Markdown) con las secciones: resumen, presupuesto de cómputo, tamaño óptimo, arquitectura, plan de paralelismo, plan de servicio y riesgos.
  \item Cuaderno o \emph{script} de Python que reproduzca todas las cifras del informe.
  \item Al menos dos diagramas: el plan de paralelismo y la arquitectura del bloque.
  \item Una tabla resumen final con $N$, $D$, $C$, GPU-horas, días, MWh, coste y tokens/s servidos.
\end{enumerate}

\paragraph{Pasos sugeridos.}
\begin{enumerate}[label=\arabic*.]
  \item Calcular $C$ a partir del clúster. Hallar el óptimo con la ley de Chinchilla (M16, §4.3) y con la regla de 20 tokens (M16, §4.5), y discutir la diferencia.
  \item Decidir $N$ y $D$ finales teniendo en cuenta el coste de inferencia (M16, §5.3), el límite de tokens únicos y el número de épocas aceptable (M16, §5.4) y la restricción de servicio.
  \item Elegir la arquitectura: $d$, $L$, cabezas, $\nkv$, $\dff$; modelo denso o MoE (M18, §1.3, §2.3 y §4.2). Comprobar que $N\approx 12Ld^{2}$ cuadra con la elección, o su equivalente MoE.
  \item Diseñar el paralelismo: TP, PP, DP, etapa de ZeRO y microlotes (M17, §5.3--§5.6). Calcular memoria por GPU, burbuja y tiempo de comunicación por paso.
  \item Diseñar el servicio: precisión de los pesos, caché KV, lote máximo y tokens/s (M19, §2.2--§2.4); decidir si se usa cuantización (M19, §4.3) o decodificación especulativa (M19, §3.4).
  \item Enumerar los riesgos (fallos de hardware, MFU inferior a la prevista, datos insuficientes, picos de pérdida) y proponer una mitigación para cada uno.
\end{enumerate}

\paragraph{Criterios de evaluación.} El proyecto se puntúa sobre 100 puntos con la rúbrica del \cref{tab:rubrica}. Se aprueba con 70 puntos o más.

\begin{table}[h]
\centering\footnotesize
\caption{Rúbrica del proyecto de fase.}\label{tab:rubrica}
\begin{tabularx}{\textwidth}{@{}>{\raggedright\arraybackslash}p{2.5cm}c>{\raggedright\arraybackslash}X>{\raggedright\arraybackslash}X>{\raggedright\arraybackslash}X@{}}
\toprule
Criterio & Peso & Excelente & Suficiente & Insuficiente \\
\midrule
Presupuesto de cómputo, tiempo, energía y coste & 20 & $C$, GPU-horas, MWh y coste correctos y reproducibles; análisis de sensibilidad a la MFU & Cifras correctas sin análisis de sensibilidad & Errores de orden de magnitud o de unidades \\
\addlinespace
Tamaño óptimo y datos & 20 & Compara Chinchilla y la regla de 20; justifica la desviación por inferencia y por datos únicos & Aplica un criterio correctamente & Elige $N$ sin justificar o ignora el límite de tokens únicos \\
\addlinespace
Arquitectura & 15 & Dimensiones coherentes con $N$; justifica GQA y la elección entre denso y MoE & Dimensiones coherentes sin discusión & Dimensiones que no cuadran con $N$ \\
\addlinespace
Plan de paralelismo & 20 & TP, PP, DP y ZeRO justificados por la topología; memoria, burbuja y comunicación calculadas & Plan viable con algún cálculo ausente & Plan que no cabe en memoria o que ignora la red \\
\addlinespace
Plan de servicio & 15 & Cumple la restricción con cálculo de caché KV y ancho de banda; valora la cuantización & Calcula los tokens/s con alguna simplificación & No comprueba la restricción \\
\addlinespace
Claridad, diagramas y reproducibilidad & 10 & Informe claro, diagramas útiles y código que reproduce todo & Informe legible y código parcial & Sin código o con cifras sin origen \\
\bottomrule
\end{tabularx}
\end{table}

\paragraph{Calificación del hito.} Nota del examen $=$ puntos$/10$; nota del proyecto $=$ puntos$/10$. Nota del hito $=0{,}6\cdot\text{examen}+0{,}4\cdot\text{proyecto}$. Para superar el hito hay que obtener al menos un 7 en cada una de las dos partes.
